% ==============================================================================
% Phân mục 08: Thảo luận, Bảng kiểm PROBAST & Tuyên bố Y đức (Discussion & Ethics)
% [PHỐI HỢP THÀNH VIÊN 2 --- BACKLOG 07 NHIỆM VỤ 2.2 & 2.3]
% ==============================================================================
\section{Thảo Luận, Thẩm Định PROBAST và Tuyên Bố Y Đức (Discussion, PROBAST \& Ethics)}
\label{sec:discussion_ethics}

\subsection{Đối chiếu Tiêu chuẩn Nguy cơ Sai lệch Quốc tế PROBAST}
Khác với 15 nghiên cứu tiền nhiệm trong tổng quan hệ thống của Hunik và cộng sự \cite{hunik2024diagnostic} đều rơi vào mức nguy cơ sai lệch cao hoặc không rõ ràng, quy trình phát triển và kiểm định M-CDSS được thiết kế ngay từ đầu để đáp ứng cả 4 miền tiêu chuẩn của công cụ \textbf{PROBAST} \cite{collins2021protocol} (Bảng~\ref{tab:probast_summary}).

\begin{table}[htbp]
\centering
\caption{Tóm tắt kết quả tự kiểm toán nguy cơ sai lệch theo 4 miền tiêu chuẩn PROBAST \cite{collins2021protocol}.}
\label{tab:probast_summary}
\footnotesize
\begin{tabularx}{\textwidth}{@{}llX@{}}
\toprule
\textbf{Miền Đánh giá (Domain)} & \textbf{Mức Sai lệch} & \textbf{Bằng chứng Thực thi Kỹ thuật trong Dự án M-CDSS} \\
\midrule
1. Đối tượng (\textit{Participants}) & \textbf{Thấp (\textit{Low})} & Đoàn hệ 2.000 hồ sơ đa phương thức cân bằng 4 nhóm bệnh; phân chia nghiêm ngặt theo định danh bệnh nhân (\textit{Strict Group Split}), kiểm toán tự động 0\% trùng lặp giữa Train/Val/Test. \\[0.4em]
2. Biến dự báo (\textit{Predictors}) & \textbf{Thấp (\textit{Low})} & Áp dụng khóa thời gian \texttt{TemporalAligner} và danh sách cấm tiêu chuẩn vàng (\texttt{[REDACTED\_DIAGNOSIS]}), đảm bảo không sử dụng thông tin phát sinh sau thời điểm chẩn đoán. \\[0.4em]
3. Kết cục (\textit{Outcome}) & \textbf{Thấp (\textit{Low})} & Nhãn chẩn đoán xác định dựa trên mã chuẩn hóa ICD-10 và đối chiếu nhất quán với tiêu chuẩn vàng trong hướng dẫn GINA 2024, GOLD 2024, ESC và Bộ Y tế. \\[0.4em]
4. Phân tích (\textit{Analysis}) & \textbf{Thấp (\textit{Low})} & Xử lý dữ liệu khuyết thiếu (WBC 20.5\%, CRP 22.4\%) bằng thuật toán MICE kèm biến chỉ báo khuyết thiếu chỉ khớp trên tập Train; hiệu chuẩn xác suất bằng Temperature Scaling (kết quả ECE thực nghiệm: \texttt{[TODO-TV2: Điền ECE thực nghiệm < 0.05]}); lượng hóa lợi ích ròng lâm sàng bằng DCA. \\
\bottomrule
\end{tabularx}
\end{table}

\begin{itemize}[leftmargin=1.8em, itemsep=0.2em]
    \item \textbf{Phụ lục Bảng Thẩm định Chi tiết 20 Tiêu chí PROBAST (Chờ bổ sung):} \texttt{[TODO-TV2 (Backlog 07 Task 2.2): Đính kèm bảng giải trình chi tiết từng câu hỏi tín hiệu (20 signaling questions) của PROBAST vào phần Supplementary Material]}.
    \item \textbf{Chứng chỉ Kiểm toán Bảo mật PHI Toàn Kho lưu trữ (Chờ bổ sung):} \texttt{[TODO-TV2 (Backlog 07 Task 2.3): Điền kết quả rà quét bảo mật PHI lần cuối trên toàn bộ repository và lịch sử git commit]}.
\end{itemize}

\subsection{Giới hạn Nghiên cứu và Tuyên bố Y đức (Limitations \& Ethical Boundaries)}
Bên cạnh những kết quả tích cực, nghiên cứu ghi nhận một số giới hạn minh bạch cần tiếp tục mở rộng trong tương lai:
\begin{itemize}[leftmargin=1.8em, itemsep=0.2em]
    \item \textbf{Đặc thù dữ liệu tổng hợp chuẩn hóa:} Mặc dù bộ dữ liệu kết hợp từ Synthea EHR và MTSamples phản ánh sát phân bố sinh lý và văn phong ghi chú y khoa, hệ thống cần được ngoại kiểm (\textit{External Validation}) thêm trên dữ liệu bệnh án điện tử thực tế đa trung tâm tại các bệnh viện tuyến đầu Việt Nam trước khi triển khai lâm sàng rộng rãi.
    \item \textbf{Phạm vi nhóm bệnh:} Hiện tại mô hình tập trung phân loại 4 hội chứng hô hấp -- tim mạch phổ biến nhất; việc mở rộng sang các bệnh lý đồng mắc phức tạp (\textit{multi-label comorbidities}) là hướng phát triển tiếp theo của kiến trúc dung hợp.
\end{itemize}

\begin{clinicalbox}{TUYÊN BỐ RANH GIỚI LÂM SÀNG VÀ ĐẠO ĐỨC Y SINH (CLINICAL DISCLAIMER)}
Hệ thống M-CDSS là một nguyên mẫu nghiên cứu khoa học phục vụ \textbf{hỗ trợ ra quyết định lâm sàng (Clinical Decision Support)}, được thiết kế để cung cấp góc nhìn tham khảo thứ hai kèm dẫn chiếu y văn minh bạch cho nhân viên y tế. Hệ thống \textbf{tuyệt đối không thay thế quyền quyết định chuyên môn của bác sĩ điều trị} và không tự động ra y lệnh trên bệnh nhân thực tế.
\end{clinicalbox}
